% GENERATED FILE. Edit canonical lessons, then run npm run generate:books.
% canonical-source-hash: fnv1a64:c048293a40bb5989
% canonical-lessons: ES-C675-afueras, ES-C675-cuadra, ES-C675-callejon, ES-C675-soportal, ES-C675-adoquin

\chapter{Las partes de la ciudad}
\label{ch:es-las-partes-de-la-ciudad}
\hlchaptermodality{675}

\begin{chapteropening}
\textbf{By the end of this chapter:} \emph{I can talk about parts of a town: the outskirts, the blocks, an alley, arcades and cobbled streets.}

Complete the last lesson of chapter 675: I can talk about parts of a town: the outskirts, the blocks, an alley, arcades and cobbled streets.
\end{chapteropening}

\section[las afueras]{las afueras — the outskirts}

\label{lesson:ES-C675-afueras}



\begin{quote}
Before the new one: say the Spanish for to travel around, then the Spanish for the route.
\end{quote}

\subsection*{You'll want to know: las afueras}
\textbf{las afueras} — \textquotedblleft{}the outskirts\textquotedblright{}.

\begin{quote}
Vivimos en las afueras de la ciudad.
\end{quote}

\begin{quote}
\emph{We live on the outskirts of the city.}
\end{quote}

Always plural: \emph{las afueras}.

\begin{cousinweb}
\emph{afueras} is built on \emph{afuera} and \emph{fuera} \textquotedblleft{}outside\textquotedblright{}, which come from Latin \emph{forās} \textquotedblleft{}out of doors\textquotedblright{} (EtymDB).
\end{cousinweb}

\begin{grammarlens}[title={What you've built}]
One more word for telling, explaining and describing at B1.
\end{grammarlens}

\subsection*{Your turn}
\begin{itemize}
\raggedright
  \item \emph{Say it:} \emph{las afueras}
  \item \emph{Say it:} \emph{las afueras}, once more
  \item \emph{Say it:} a sentence of your own with \emph{las afueras}
  \item \emph{Recall:} say the Spanish for to travel around, then the Spanish for the route, then say \emph{las afueras} again
\end{itemize}

\begin{tcolorbox}[breakable,colback=teal!4,colframe=teal!35!black,title={Before you move on}]
What does las afueras mean? (\textquotedblleft{}The outskirts\textquotedblright{}.) Say it once more.
\end{tcolorbox}
\section[la cuadra]{la cuadra — the block (of a street) (Latin America); the stable}

\label{lesson:ES-C675-cuadra}



\begin{quote}
Before the new one: say the Spanish for the route, then the Spanish for the outskirts.
\end{quote}

\subsection*{You'll want to know: la cuadra}
\textbf{la cuadra} — \textquotedblleft{}the block (of a street) (Latin America); the stable\textquotedblright{}.

\begin{quote}
El hotel está a dos cuadras.
\end{quote}

\begin{quote}
\emph{The hotel is two blocks away.}
\end{quote}

Spain says \emph{la manzana} for a block; there \emph{la cuadra} is a stable.

\begin{cousinweb}
\emph{cuadra} comes from Latin \emph{quadra} \textquotedblleft{}a square\textquotedblright{} (Lewis \& Short): a city block is a square of streets.
\end{cousinweb}

\begin{grammarlens}[title={What you've built}]
One more word for telling, explaining and describing at B1.
\end{grammarlens}

\subsection*{Your turn}
\begin{itemize}
\raggedright
  \item \emph{Say it:} \emph{la cuadra}
  \item \emph{Say it:} \emph{la cuadra}, once more
  \item \emph{Say it:} a sentence of your own with \emph{la cuadra}
  \item \emph{Recall:} say the Spanish for the route, then the Spanish for the outskirts, then say \emph{la cuadra} again
\end{itemize}

\begin{tcolorbox}[breakable,colback=teal!4,colframe=teal!35!black,title={Before you move on}]
What does la cuadra mean? (\textquotedblleft{}The block (of a street) (Latin America); the stable\textquotedblright{}.) Say it once more.
\end{tcolorbox}
\section[el callejón]{el callejón — the alley, the lane}

\label{lesson:ES-C675-callejon}



\begin{quote}
Before the new one: say the Spanish for the outskirts, then the Spanish for the block.
\end{quote}

\subsection*{You'll want to know: el callejón}
\textbf{el callejón} — \textquotedblleft{}the alley, the lane\textquotedblright{}.

\begin{quote}
El bar está en un callejón.
\end{quote}

\begin{quote}
\emph{The bar is down an alley.}
\end{quote}

\emph{Un callejón sin salida} — a dead end.

\begin{cousinweb}
\emph{callejón} is built on \emph{calleja} \textquotedblleft{}a little street\textquotedblright{} (EtymDB), and \emph{calle} is inherited from Latin \emph{callis} \textquotedblleft{}a footpath\textquotedblright{}.
\end{cousinweb}

\begin{grammarlens}[title={What you've built}]
One more word for telling, explaining and describing at B1.
\end{grammarlens}

\subsection*{Your turn}
\begin{itemize}
\raggedright
  \item \emph{Say it:} \emph{el callejón}
  \item \emph{Say it:} \emph{el callejón}, once more
  \item \emph{Say it:} a sentence of your own with \emph{el callejón}
  \item \emph{Recall:} say the Spanish for the outskirts, then the Spanish for the block, then say \emph{el callejón} again
\end{itemize}

\begin{tcolorbox}[breakable,colback=teal!4,colframe=teal!35!black,title={Before you move on}]
What does el callejón mean? (\textquotedblleft{}The alley, the lane\textquotedblright{}.) Say it once more.
\end{tcolorbox}
\section[el soportal]{el soportal — the arcade (covered walkway)}

\label{lesson:ES-C675-soportal}



\begin{quote}
Before the new one: say the Spanish for the block, then the Spanish for the alley.
\end{quote}

\subsection*{You'll want to know: el soportal}
\textbf{el soportal} — \textquotedblleft{}the arcade (covered walkway)\textquotedblright{}.

\begin{quote}
Con la lluvia caminamos bajo los soportales.
\end{quote}

\begin{quote}
\emph{In the rain we walk under the arcades.}
\end{quote}

Often plural: \emph{los soportales} of a main square.

\begin{cousinweb}
\emph{soportal} is built from \emph{so-} \textquotedblleft{}under\textquotedblright{}, from Latin \emph{sub}, and \emph{portal} \textquotedblleft{}doorway\textquotedblright{}, from Latin \emph{porta} \textquotedblleft{}gate\textquotedblright{}.
\end{cousinweb}

\begin{grammarlens}[title={What you've built}]
One more word for telling, explaining and describing at B1.
\end{grammarlens}

\subsection*{Your turn}
\begin{itemize}
\raggedright
  \item \emph{Say it:} \emph{el soportal}
  \item \emph{Say it:} \emph{el soportal}, once more
  \item \emph{Say it:} a sentence of your own with \emph{el soportal}
  \item \emph{Recall:} say the Spanish for the block, then the Spanish for the alley, then say \emph{el soportal} again
\end{itemize}

\begin{tcolorbox}[breakable,colback=teal!4,colframe=teal!35!black,title={Before you move on}]
What does el soportal mean? (\textquotedblleft{}The arcade (covered walkway)\textquotedblright{}.) Say it once more.
\end{tcolorbox}
\section[el adoquín]{el adoquín — the paving stone, the cobblestone}

\label{lesson:ES-C675-adoquin}



\begin{quote}
Before the new one: say the Spanish for the alley, then the Spanish for the arcade.
\end{quote}

\subsection*{You'll want to know: el adoquín}
\textbf{el adoquín} — \textquotedblleft{}the paving stone, the cobblestone\textquotedblright{}.

\begin{quote}
Las calles del centro son de adoquines.
\end{quote}

\begin{quote}
\emph{The streets in the centre are cobbled.}
\end{quote}

Hard work for a suitcase on wheels.

\begin{cousinweb}
\emph{adoquín} comes from Andalusian Arabic \emph{ad-dukkán} \textquotedblleft{}stone bench\textquotedblright{}, from Arabic \emph{dukkān} (EtymDB).
\end{cousinweb}

\begin{grammarlens}[title={What you've built}]
One more word for telling, explaining and describing at B1.
\end{grammarlens}

\subsection*{Your turn}
\begin{itemize}
\raggedright
  \item \emph{Say it:} \emph{el adoquín}
  \item \emph{Say it:} \emph{el adoquín}, once more
  \item \emph{Say it:} a sentence of your own with \emph{el adoquín}
  \item \emph{Recall:} say the Spanish for the alley, then the Spanish for the arcade, then say \emph{el adoquín} again
\end{itemize}

\begin{tcolorbox}[breakable,colback=teal!4,colframe=teal!35!black,title={Before you move on}]
What does el adoquín mean? (\textquotedblleft{}The paving stone, the cobblestone\textquotedblright{}.) Say it once more.
\end{tcolorbox}
